% =====================================================================
% Elsevier elsarticle 期刊投稿最小可编译骨架
% 编译: latexmk -pdf elsevier_elsarticle.tex      (pdfLaTeX)
% 投稿审稿版: [preprint,review,12pt] -> 出行号 + 宽行距便于批注
% 终排版式:   [final,5p,times,twocolumn]  (1p单栏/3p/5p)
% 参考文献模型按目标期刊 guide for authors 指定:
%   elsarticle-num (数字) / elsarticle-num-names / elsarticle-harv (著者-年)
% =====================================================================
\documentclass[preprint,review,12pt]{elsarticle}

\usepackage{amsmath}
\usepackage{graphicx}
\usepackage{lineno}   % review 选项会启用行号

\journal{Journal Name}

\begin{document}

\begin{frontmatter}

\title{Paper Title Here}

\author[inst1]{First Author\corref{cor1}}
\ead{first@example.com}
\cortext[cor1]{Corresponding author}

\author[inst2]{Second Author}
\ead{second@example.com}

\affiliation[inst1]{organization={First Institution},
            city={City}, country={Country}}
\affiliation[inst2]{organization={Second Institution},
            city={City}, country={Country}}

\begin{abstract}
This is the abstract. State problem, method, results and contributions in one paragraph.
\end{abstract}

\begin{keyword}
keyword1 \sep keyword2 \sep keyword3
\end{keyword}

\end{frontmatter}

\section{Introduction}
Body text. Citation example~\cite{ref1}. See Fig.~\ref{fig:arch} and Eq.~\eqref{eq:main}.

\section{Method}
\begin{equation}\label{eq:main}
  y = f(x) + \epsilon
\end{equation}

\begin{figure}[t]
  \centering
  % \includegraphics[width=0.6\textwidth]{fig1.pdf}
  \rule{0.6\textwidth}{3cm}   % 占位框, 有图时替换
  \caption{System architecture.}
  \label{fig:arch}
\end{figure}

\section{Conclusion}
Summary.

% --- 参考文献: 按期刊指定模型 -------------------------------------
\bibliographystyle{elsarticle-num}
% \bibliography{refs}
% 演示用内联:
\begin{thebibliography}{1}
\bibitem{ref1}
A.~Author, Title of paper, Journal 1 (1) (2024) 1--10.
\end{thebibliography}

\end{document}
